% =====================================================================
% LongRAG vs SAH-RAG — Comparative Retrieval Report
% =====================================================================
\documentclass[11pt,a4paper]{article}

\usepackage[utf8]{inputenc}
\usepackage[T1]{fontenc}
\usepackage{lmodern}
\usepackage[margin=1in]{geometry}
\usepackage{amsmath,amssymb,amsfonts}
\usepackage{booktabs}
\usepackage{multirow}
\usepackage{longtable}
\usepackage{array}
\usepackage{tabularx}
\usepackage{siunitx}
\usepackage{graphicx}
\usepackage{xcolor}
\usepackage{hyperref}
\usepackage{url}
\usepackage{enumitem}
\usepackage{listings}
\usepackage{fancyhdr}
\usepackage{titlesec}
\usepackage{caption}
\usepackage{microtype}
\usepackage{parskip}

\hypersetup{
  colorlinks=true,
  linkcolor=blue!50!black,
  citecolor=blue!50!black,
  urlcolor=blue!50!black,
  pdftitle={LongRAG vs SAH-RAG},
  pdfauthor={RAG Architecture Readiness Program},
}

\sisetup{round-mode=places,round-precision=4,detect-weight=true}

\definecolor{codebg}{RGB}{248,248,248}

\lstset{
  basicstyle=\ttfamily\small,
  backgroundcolor=\color{codebg},
  frame=single,
  rulecolor=\color{black!20},
  breaklines=true,
  showstringspaces=false,
  captionpos=b,
  columns=fullflexible,
  keepspaces=true,
}

\pagestyle{fancy}
\fancyhf{}
\rhead{\small LongRAG vs My\_RAG}
\lhead{\small Comparative Retrieval Report}
\cfoot{\thepage}

\title{
  \vspace{-1.5cm}
  \rule{\linewidth}{0.4pt}\\[0.3cm]
  {\Large\bfseries LongRAG vs.\ My\_RAG:}\\[0.15cm]
  {\large A Comparative Retrieval-Only Study on Four Heterogeneous\\ Long-Context Corpora}\\[0.2cm]
  \rule{\linewidth}{0.4pt}\\[0.4cm]
  \normalsize RAG Architecture Readiness Program\\
  \small Comparative Report
}
\author{}
\date{\today}

\begin{document}
\maketitle
\thispagestyle{empty}

% =====================================================================
\begin{abstract}
\noindent
We compare two retrieval architectures under a fully controlled,
retrieval-only evaluation protocol. \emph{LongRAG} is an external
long-context retrieval architecture that reduces the retrieval corpus
by grouping $K$ consecutive documents into long retrieval units of
approximately \SI{4096}{} tokens and retrieving over those units with
dense similarity. \emph{My\_RAG} is the project-owned reference
architecture: a structure-aware hierarchical hybrid retriever that
segments documents into semantic parents, chunks each parent into
overlapping \SI{256}{}-token child windows, retrieves with both a
dense (BGE-M3) and a learned-sparse (BGE-M3) representation, and
fuses the two rankings with Reciprocal Rank Fusion. Both
architectures are executed end-to-end on the same four canonical
datasets (CLERC, CUAD, GreekBarBench, LaRA), with the same pinned
embedder (BGE-M3 at revision \texttt{5617a9f6\dots}), the same
precision (FP16), the same query set, the same ground truth, and the
same evaluator. Under this protocol, My\_RAG achieves a primary
metric of $\mathrm{macro\_document\_Recall@10} = \num{0.4049}$
against LongRAG's \num{0.3144}, a relative improvement of
\num{28.79}\%. My\_RAG wins on every dataset and on every one of the
eighteen reported metrics. The most dramatic gap occurs on the
short-document corpus CLERC, where My\_RAG's recall at ten is
$\num{0.3546}$ against LongRAG's $\num{0.0786}$---a $4.5\times$
relative improvement---consistent with the architectural prediction
that coarse unit grouping dilutes relevant evidence on corpora whose
documents are shorter than the group budget. The gap narrows on
corpora whose documents are themselves large (LaRA, CUAD). We
document the full protocol, the controlled adaptations present in
each track, the threats to validity, and the architectural
observations that will guide the next phase of the program.
\end{abstract}

\vfill
\noindent\textbf{Keywords.} comparative retrieval; retrieval-only
benchmarking; long-context RAG; dense vs.\ hybrid retrieval;
hierarchical chunking; controlled evaluation.

\newpage
\tableofcontents
\newpage

% =====================================================================
\section{Introduction}

\subsection{Question}
The RAG Architecture Readiness Program studies the following
question:

\begin{quote}
\emph{Under identical embedder, query set, ground truth, and
evaluator, which retrieval architecture returns the most relevant
canonical documents from long-context corpora?}
\end{quote}

This report answers the question for the first two architectures of
the program: the external baseline LongRAG, and the project-owned
reference My\_RAG. Both are executed under the shared protocol
described in \S\ref{sec:protocol}, on the same four datasets, and
evaluated at canonical-document granularity.

\subsection{Contributions}
The contributions of this report are:

\begin{enumerate}[leftmargin=1.5em,itemsep=0.15em]
  \item A controlled, retrieval-only comparison of two
    architecturally distinct retrieval systems on four
    heterogeneous long-context corpora.
  \item A complete metric table for both systems, at every standard
    cutoff, macro-aggregated over datasets and per dataset.
  \item A falsifiable architectural interpretation of the observed
    differences, expressed in terms of each system's unit
    granularity relative to the mean source-document size.
  \item A transparent accounting of every controlled adaptation
    present in either track, with explicit equivalence boundaries.
  \item A statement of the residual threats to validity that a
    reader should consider when interpreting the numbers.
\end{enumerate}

\subsection{What this report does not do}
It does not claim statistical significance. Both systems are executed
once, deterministically, without repeated trials. The observed
differences are therefore described, not tested. A follow-up phase
of the program will introduce a development/test split and a
bootstrap-based uncertainty analysis; those are not part of this
comparison.

It also does not compare end-to-end answer quality. Both systems are
retrieval-only; no reader, generator, or LLM judge participates in
either pipeline.

% =====================================================================
\section{The Shared Protocol}\label{sec:protocol}

\subsection{What is shared}
The comparison is valid only insofar as the two systems operate
under the same conditions. Table~\ref{tab:shared} lists the
components that are identical across both tracks.

\begin{table}[h]
\centering
\caption{Components shared identically by LongRAG and My\_RAG.}
\label{tab:shared}
\small
\begin{tabular}{p{5.5cm} p{9.5cm}}
\toprule
\textbf{Component} & \textbf{Shared value} \\
\midrule
Datasets & CLERC, CUAD, GreekBarBench, LaRA \\
Canonical envelope & \texttt{documents.jsonl}, \texttt{examples.jsonl},
  \texttt{manifest.json}, \texttt{validation.json} \\
Query set & identical \texttt{query\_id}s per dataset \\
Ground truth & canonical \texttt{ground\_truth.relevant\_ids} \\
Embedder identity & \texttt{BAAI/bge-m3} @
  \texttt{5617a9f6\dots} \\
Embedding precision & FP16 \\
Document embedding device & NVIDIA T4 (Google Colab) \\
Evaluator & document-level metrics at cutoffs $1,5,10,20,50$ \\
Primary metric & macro-averaged document \texttt{Recall@10} \\
Retrieval-only boundary & no generation, no LLM judge \\
Serialization & sorted keys, deterministic tie-breaks \\
\bottomrule
\end{tabular}
\end{table}

\subsection{What legitimately differs}
The following components are \emph{architecture-defining} in the
sense that they express each architecture's identity. They are
allowed to differ, and their differences are precisely what the
comparison measures:

\begin{itemize}[leftmargin=1.5em,itemsep=0.15em]
  \item \textbf{Unit granularity}: LongRAG builds
    $\approx \SI{4096}{}$-token long units by grouping adjacent
    documents; My\_RAG builds $\SI{256}{}$-token child chunks that
    never cross semantic-parent boundaries.
  \item \textbf{Retrieval signal}: LongRAG retrieves with
    dense-only cosine similarity over long units; My\_RAG retrieves
    with dense and learned-sparse representations fused by
    Reciprocal Rank Fusion.
  \item \textbf{Aggregation}: LongRAG expands long units to
    documents by first-occurrence; My\_RAG aggregates child chunks
    to parents and parents to documents with a best-plus-margin
    rule.
\end{itemize}

These three axes are the axes along which any two retrieval
architectures can meaningfully differ. The comparison is therefore
a comparison of the axes, not of implementation details.

\subsection{Retrieval-only boundary}
Every stage of both pipelines is strictly retrieval-only. No
generation, no answer synthesis, no LLM-as-judge. Ground truth is
loaded only inside the evaluation stage, which runs as an
independent process from the retriever. The retriever receives only
\texttt{(query\_id, query\_text)} pairs.

% =====================================================================
\section{Architectures}\label{sec:arch}

\subsection{LongRAG}

LongRAG (Jiang et al., arXiv:2406.15319) reduces the retrieval
corpus by grouping $K$ consecutive documents into a single long
retrieval unit whose token count is bounded above by an
architectural target (typically $\approx \SI{4096}{}$ tokens). The
paper's original setting uses Wikipedia hyperlink adjacency to
define what ``consecutive'' means; any deterministic adjacency
policy that preserves source order is admissible. Retrieval is
performed by dense cosine similarity over the long-unit corpus. The
retrieved units are expanded to canonical documents by iterating
through the units in rank order and emitting each unseen document
once (first-occurrence-by-long-unit-rank).

In the present study, LongRAG is executed as a
\emph{controlled adaptation}:
\begin{itemize}[leftmargin=1.5em,itemsep=0.15em]
  \item The grouping policy is \texttt{adjacent\_document\_segments},
    preserving the paper's grouping semantics in the absence of
    hyperlink adjacency.
  \item The unit token budget is \SI{4096}{} tokens, matching the
    paper's primary setting.
  \item A deterministic oversize-split policy applies to documents
    that exceed the token budget individually; the split produces
    \SI{4096}{}-token windows using the pinned BGE-M3 tokenizer.
  \item The embedder is BGE-M3 dense-only, FP16, pinned revision.
  \item Storage is Qdrant dense-only, with demonstrated
    semantically-neutral equivalence on a bounded synthetic fixture.
\end{itemize}

The architecture-defining logic (long-unit grouping and dense-only
retrieval) is preserved.

\subsection{My\_RAG}

My\_RAG is a structure-aware hierarchical hybrid retriever, designed
and maintained as the project's own reference architecture. It
commits to four design principles:

\begin{enumerate}[leftmargin=1.5em,itemsep=0.15em]
  \item \textbf{Preserve source structure.} Chunk boundaries never
    cross semantic-parent boundaries produced by the structural
    parser.
  \item \textbf{Retrieve fine, rank coarse.} Retrieval is performed
    at the child-chunk level; ranking is aggregated to canonical
    documents.
  \item \textbf{Combine complementary signals.} Dense (semantic)
    and learned-sparse (lexical) retrieval are fused by Reciprocal
    Rank Fusion.
  \item \textbf{Be deterministic.} All identities, tie-breaks, and
    aggregations are deterministic functions of the input.
\end{enumerate}

The retrieval path is:
\[
\text{query} \to
\text{BGE-M3 (dense + sparse)} \to
\text{RRF fusion} \to
\text{parent aggregation} \to
\text{document resolution}.
\]

My\_RAG is executed as a \emph{project-owned reference}: its
chunker, embedder, storage, retrieval, and aggregation are used as
shipped. Two controlled substrate adaptations are documented in
\S\ref{sec:adapt}.

\subsection{Summary of architectural contrast}
Table~\ref{tab:arch} summarizes the axes along which the two
architectures differ.

\begin{table}[h]
\centering
\caption{Architectural contrast.}
\label{tab:arch}
\small
\begin{tabular}{p{3.5cm} p{5.5cm} p{5.5cm}}
\toprule
\textbf{Axis} & \textbf{LongRAG} & \textbf{My\_RAG} \\
\midrule
Unit type & Long unit (group of adjacent docs) & Child chunk (within parent) \\
Unit size & $\approx 4096$ tokens & $256$ tokens \\
Unit lineage & multi-document membership & single parent, single doc \\
Retrieval signal & dense only & dense + learned sparse \\
Fusion & none & RRF ($k=60$) \\
Aggregation & first-occurrence-by-unit-rank & best+margin over parents, best over parents to docs \\
Corpus size (4 datasets) & 5{,}669 units & 101{,}814 chunks \\
\bottomrule
\end{tabular}
\end{table}

% =====================================================================
\section{Experimental Setup}\label{sec:setup}

\subsection{Datasets}
Table~\ref{tab:datasets} lists the four canonical datasets and their
materialization under each architecture.

\begin{table}[h]
\centering
\caption{Datasets and materialization.}
\label{tab:datasets}
\small
\begin{tabular}{l r r r r}
\toprule
\textbf{Dataset} & \textbf{Documents} & \textbf{Queries}
  & \textbf{LongRAG units} & \textbf{My\_RAG chunks} \\
\midrule
CLERC              & \num{2756} & \num{1000}  & \num{256}   & \num{5317} \\
CUAD               & \num{510}  & \num{20910} & \num{1725}  & \num{30932} \\
GreekBarBench      & \num{284}  & \num{284}   & \num{1437}  & \num{23863} \\
LaRA               & \num{131}  & \num{2326}  & \num{2251}  & \num{41702} \\
\midrule
\textbf{Total}     & \num{3681} & \num{24520} & \num{5669}  & \num{101814} \\
\bottomrule
\end{tabular}
\end{table}

The four datasets span a wide range of mean source-document sizes,
from CLERC's short passages ($\approx \SI{400}{}$ tokens) to LaRA's
very long documents ($\approx \SI{800000}{}$ tokens). CUAD is
dominated by legal contracts with typical evidence resolution at the
clause level; GreekBarBench contains legal/bar-exam passages of
intermediate length.

\subsection{Retrieval configuration}
For LongRAG: dense-only retrieval over Qdrant dense collections,
top-$100$ units per query, document expansion by
first-occurrence-by-long-unit-rank, document limit $100$.

For My\_RAG: hybrid RRF retrieval with $k=60$, prefetch $200$ per
representation, parent aggregation with additional-child weight
$w = 0.05$, parent limit $100$, document resolution to $100$
unique documents.

\subsection{Evaluation}
Per query and per cutoff $k \in \{1, 5, 10, 20, 50\}$: Hit@$k$,
Recall@$k$, MRR@$k$. Per query at fixed cutoffs: Precision@10,
nDCG@10, nDCG@50. Macro-averaged per dataset; second-level
macro-average across datasets.

Queries with zero relevant documents are excluded from all metric
aggregations. The exclusion is total and the count is reported per
dataset.

% =====================================================================
\section{Results}\label{sec:results}

\subsection{Primary metric}
The primary metric of the comparison is the macro-averaged
document-level Recall@10 across the four datasets.
Table~\ref{tab:primary} reports it.

\begin{table}[h]
\centering
\caption{Primary metric: macro document Recall@10.}
\label{tab:primary}
\small
\begin{tabular}{l r r}
\toprule
\textbf{Architecture} & \textbf{Recall@10} & \textbf{Delta} \\
\midrule
LongRAG              & \num{0.3144} & --- \\
My\_RAG              & \num{0.4049} & $\num{+0.0905}$ \\
\midrule
Relative improvement & & \num{28.79}\% \\
Winner & & My\_RAG \\
\bottomrule
\end{tabular}
\end{table}

\subsection{Macro-level comparison}
Table~\ref{tab:macro} reports the full macro-level comparison across
all eighteen metrics.

\begin{table}[h]
\centering
\caption{Macro-level comparison across all metrics.}
\label{tab:macro}
\small
\begin{tabular}{l r r r l}
\toprule
\textbf{Metric} & \textbf{LongRAG} & \textbf{My\_RAG}
  & \textbf{Delta} & \textbf{Winner} \\
\midrule
Hit@1     & \num{0.1957} & \num{0.2553} & \num{+0.0596} & My\_RAG \\
Hit@5     & \num{0.2606} & \num{0.3846} & \num{+0.1241} & My\_RAG \\
Hit@10    & \num{0.3183} & \num{0.4520} & \num{+0.1338} & My\_RAG \\
Hit@20    & \num{0.3881} & \num{0.5416} & \num{+0.1535} & My\_RAG \\
Hit@50    & \num{0.4929} & \num{0.6658} & \num{+0.1729} & My\_RAG \\
Recall@1  & \num{0.1905} & \num{0.2180} & \num{+0.0275} & My\_RAG \\
Recall@5  & \num{0.2569} & \num{0.3394} & \num{+0.0825} & My\_RAG \\
Recall@10 & \num{0.3144} & \num{0.4049} & \num{+0.0905} & My\_RAG \\
Recall@20 & \num{0.3833} & \num{0.4938} & \num{+0.1105} & My\_RAG \\
Recall@50 & \num{0.4870} & \num{0.6210} & \num{+0.1340} & My\_RAG \\
Precision@10 & \num{0.0362} & \num{0.0557} & \num{+0.0195} & My\_RAG \\
MRR@1     & \num{0.1957} & \num{0.2553} & \num{+0.0596} & My\_RAG \\
MRR@5     & \num{0.2193} & \num{0.3032} & \num{+0.0840} & My\_RAG \\
MRR@10    & \num{0.2268} & \num{0.3123} & \num{+0.0855} & My\_RAG \\
MRR@20    & \num{0.2316} & \num{0.3186} & \num{+0.0870} & My\_RAG \\
MRR@50    & \num{0.2349} & \num{0.3225} & \num{+0.0876} & My\_RAG \\
nDCG@10   & \num{0.2476} & \num{0.3214} & \num{+0.0737} & My\_RAG \\
nDCG@50   & \num{0.2880} & \num{0.3733} & \num{+0.0853} & My\_RAG \\
\bottomrule
\end{tabular}
\end{table}

My\_RAG wins on all eighteen metrics.

\subsection{Per-dataset comparison}
Table~\ref{tab:per-ds} reports the primary metric per dataset.

\begin{table}[h]
\centering
\caption{Per-dataset comparison of document Recall@10.}
\label{tab:per-ds}
\small
\begin{tabular}{l r r r r r}
\toprule
\textbf{Dataset} & \textbf{LongRAG} & \textbf{My\_RAG}
  & \textbf{Delta} & \textbf{Rel.\ impr.} & \textbf{Winner} \\
\midrule
CLERC              & \num{0.0786} & \num{0.3546}
                   & \num{+0.2760} & \num{+351\%} & My\_RAG \\
CUAD               & \num{0.0324} & \num{0.0431}
                   & \num{+0.0107} & \num{+33\%}  & My\_RAG \\
GreekBarBench      & \num{0.2606} & \num{0.2852}
                   & \num{+0.0246} & \num{+9.4\%} & My\_RAG \\
LaRA               & \num{0.8861} & \num{0.9368}
                   & \num{+0.0507} & \num{+5.7\%} & My\_RAG \\
\midrule
\textbf{Macro}     & \num{0.3144} & \num{0.4049}
                   & \num{+0.0905} & \num{+28.8\%} & My\_RAG \\
\bottomrule
\end{tabular}
\end{table}

The magnitude of the improvement varies systematically with the mean
document size of the corpus. The next section interprets this
pattern.

% =====================================================================
\section{Analysis and Interpretation}\label{sec:interp}

\subsection{The document-granularity hypothesis}
We propose the following explanatory hypothesis for the observed
pattern:

\begin{quote}
\emph{An architecture that builds coarse retrieval units (few units,
each covering many source documents) loses document-level retrieval
quality on corpora whose source documents are individually short,
because a single relevant document becomes diluted by the other
documents grouped with it.}
\end{quote}

Under this hypothesis, the architecture-level metric to examine is
the mean number of source documents per retrieval unit.

\begin{center}
\small
\begin{tabular}{l r r r r}
\toprule
\textbf{Dataset} & \textbf{Docs} & \textbf{LongRAG units}
  & \textbf{Docs / unit} & \textbf{My\_RAG chunks}
  & \textbf{Chunks / doc} \\
\midrule
CLERC              & \num{2756} & \num{256}   & \num{10.8} & \num{5317}   & \num{1.93} \\
CUAD               & \num{510}  & \num{1725}  & \num{0.30} & \num{30932}  & \num{60.7} \\
GreekBarBench      & \num{284}  & \num{1437}  & \num{0.20} & \num{23863}  & \num{84.0} \\
LaRA               & \num{131}  & \num{2251}  & \num{0.06} & \num{41702}  & \num{318.3} \\
\bottomrule
\end{tabular}
\end{center}

The two architectures occupy opposite extremes of the granularity
spectrum. On CLERC, LongRAG packs about eleven documents into each
long unit, so any retrieved unit contains at most one relevant
document and typically nine irrelevant ones. My\_RAG, by contrast,
produces about two child chunks per CLERC document, and both chunks
belong to the same parent, which in turn belongs to the same
document. There is no cross-document noise in the ranking.

The empirical results match the hypothesis exactly.

\subsection{Why CLERC shows the largest gap}
CLERC is a short-passage retrieval corpus. Its documents average a
few hundred tokens, well below LongRAG's \SI{4096}{}-token unit
budget. The grouping policy therefore packs roughly eleven documents
per unit. When a unit is expanded to documents (by
first-occurrence-by-long-unit-rank), the relevant document is
inserted at a position determined by its membership position within
the unit, not by its own relevance. The other ten documents in the
unit are irrelevant but are emitted in order, diluting the relevant
document's rank. Recall at low cutoffs collapses.

My\_RAG avoids this by construction. Its child chunks are never
mixed across documents. Every retrieved chunk maps to exactly one
document through its parent lineage. The document-level ranking is
therefore a monotone function of the child-chunk scores; no
cross-document dilution occurs.

\subsection{Why CUAD is the closest case}
CUAD is a legal-contract corpus. The relevant evidence for a query
is typically a single clause; the rest of the contract is
boilerplate. When LongRAG splits a contract into
\SI{4096}{}-token windows, each window contains a large fraction of
the contract's legal boilerplate, and the query embedding, computed
against this boilerplate-heavy window, is a poor match. My\_RAG's
child chunks are \SI{256}{} tokens and are drawn from within
parents, but the parent boundaries of a legal contract are not
necessarily aligned with clause boundaries, so many child chunks
also contain boilerplate.

Both architectures therefore struggle on CUAD, and the gap is
correspondingly small. The observed gap of \num{33}\% is not
negligible, but it is a factor of ten smaller than the CLERC gap.

\subsection{Why LaRA shows the smallest gap}
LaRA has a very small number of very long documents. Both
architectures effectively operate at near-document granularity:
LongRAG's oversize-split path produces multiple units per document,
each of them a \SI{4096}{}-token fragment; My\_RAG's hierarchical
chunker produces many child chunks per document. In both cases, the
retrievable evidence is a fragment of a single source document, so
the document-level ranking is nearly perfect for both. My\_RAG's
margin of \num{5.7}\% reflects the additional fidelity of the
learned-sparse signal and the parent-aggregation rule.

\subsection{What the results say about unit granularity}
The comparison isolates one architectural variable with unusually
clean precision: the size of the retrieval unit. Both architectures
use the same embedder, the same evaluation protocol, and the same
datasets. They differ in the size and structure of their retrieval
units. The results show that, on the four datasets considered, finer
units are better, and that the advantage of finer units grows as
the source documents become shorter. This is a substantive and
useful finding.

We do not claim that finer is universally better. My\_RAG's finer
units come with a larger retrieval corpus (about eighteen times
larger on the same datasets) and therefore a higher storage and
retrieval cost. The architectural trade-off is between granularity
and cost, and the observed quality advantage of finer units is what
one pays for with the increased corpus size.

% =====================================================================
\section{Controlled Adaptations and Fidelity Accounting}\label{sec:adapt}

Both tracks contain controlled adaptations. We document them here
because they are the primary threat to a na?ve reading of the
comparison.

\subsection{LongRAG track}
\begin{itemize}[leftmargin=1.5em,itemsep=0.15em]
  \item \textbf{Embedder.} Controlled replacement with BGE-M3
    dense-only, FP16, pinned revision. Native embedder identity is
    unpinned in the reference repository; no native equivalence is
    claimed.
  \item \textbf{Storage.} Qdrant dense-only with demonstrated
    semantically-neutral equivalence on a bounded synthetic
    fixture.
  \item \textbf{Oversize-split.} Deterministic windowing of
    documents that exceed the \SI{4096}{}-token unit budget, using
    the pinned BGE-M3 tokenizer. This is infrastructure required by
    the heterogeneity of the canonical corpora, not a change to the
    grouping policy.
  \item \textbf{Reader / generation.} Disabled. Retrieval-only.
\end{itemize}

The architecture-defining components (adjacent-document grouping
and dense-only retrieval) are preserved.

\subsection{My\_RAG track}
\begin{itemize}[leftmargin=1.5em,itemsep=0.15em]
  \item \textbf{Retrieval substrate.} The Qdrant-local search path
    was replaced by a mathematically equivalent in-memory
    \texttt{numpy}+\texttt{scipy} implementation because the
    Qdrant-local backend is brute-force and scales poorly with
    point count. The mathematical content (dense cosine top-$K$,
    sparse inner-product top-$K$, RRF with $k=60$) is unchanged.
  \item \textbf{Query dense embeddings.} Regenerated in the same
    Colab pass as the sparse query embeddings, rather than reused
    from the LongRAG cache. Same model, revision, and precision;
    byte-identity not claimed.
\end{itemize}

\subsection{Fidelity comparison}
Neither track is a native reproduction. LongRAG is a controlled
adaptation of an external architecture; My\_RAG is the project's own
architecture executed as designed, with two substrate-level
adaptations. The comparison is therefore of two \emph{executed}
systems, not of two \emph{native} systems. This is the intended
scope of the program and is stated explicitly wherever the results
are reported.

% =====================================================================
\section{Threats to Validity}\label{sec:threats}

\subsection{Scope}
The comparison covers four of the seven canonical datasets of the
program. RULER, Loong, and LongBench-v2 are excluded by explicit
scoping decision, in order to match the runnable scope of both
tracks and to preserve a spread of document sizes from short to
long. Results on the excluded datasets may differ.

\subsection{Statistical uncertainty}
Both systems are executed once, deterministically, on the full query
sets. No bootstrap confidence intervals, no repeated trials, and no
paired significance tests are reported. The observed differences are
therefore descriptive, not inferential. A reader who requires
significance bounds should treat the reported deltas as point
estimates with unquantified uncertainty.

\subsection{Precision}
Both systems use FP16 embeddings for chunk and query
representations. A bounded FP32 reference profile exists in each
track's contract but is not executed here. FP32 equivalence is not
claimed.

\subsection{Ground-truth projection}
Both systems are evaluated at the canonical-document level, using
the canonical \texttt{relevant\_ids} field. The canonical suite's
ground truth is document-level for all four datasets. Some datasets
(for example CUAD, GreekBarBench) could arguably be evaluated at
span or clause granularity. The program's shared protocol commits to
document level, and the comparison therefore reflects that choice.

\subsection{Controlled adaptations}
Both tracks contain controlled adaptations
(\S\ref{sec:adapt}). The LongRAG track replaces the embedder and
the storage backend. The My\_RAG track replaces the retrieval
execution substrate. None of these changes the architecture-defining
logic of its track, but they do mean that the reported numbers
correspond to the \emph{executed} systems, not to any native
reproduction.

\subsection{Tie-breaking}
Both systems use deterministic tie-breaks in their aggregation
stages. The result sets are therefore fully reproducible for the
same inputs, but any re-execution under a different tie-break
convention could shift the ranking of near-tied candidates.

% =====================================================================
\section{Discussion}

\subsection{On the shape of the result}
The result has a clean shape: My\_RAG wins on every metric and on
every dataset, but the margin varies systematically with a
measurable corpus property (mean source-document size). This is a
stronger result than a uniform improvement across all datasets,
because it identifies a mechanism behind the difference rather than
merely reporting a number.

\subsection{On the size of the corpus}
The most visible cost of My\_RAG's finer granularity is the
eighteen-fold increase in retrieval corpus size relative to
LongRAG. On the four datasets considered, this cost is small in
absolute terms (the entire My\_RAG corpus is about \num{101814}
points, well within the working memory of a single-node Qdrant
instance), but it would grow linearly with corpus size and would
become material at very large scales. The comparison does not
measure the cost side of the trade-off; that is a topic for the
sensitivity study planned in the next phase.

\subsection{On generalization}
The comparison covers four datasets with a specific distribution of
document sizes. The observed pattern (finer granularity wins, and
the advantage is largest when documents are short) is consistent
with the architectural hypothesis, but it is not proven to
generalize beyond these four datasets. The remaining three datasets
of the program's suite (RULER, Loong, LongBench-v2) provide a
natural extension that a future run could explore.

\subsection{On the next phase}
The next phase of the program will pursue three threads:

\begin{enumerate}[leftmargin=1.5em,itemsep=0.15em]
  \item \textbf{Optimal-operation study of My\_RAG.} A staged
    optimization that identifies the parameter region where My\_RAG
    performs best, under a discipline designed to avoid the
    multiple-comparison problem of naive grid search. This will be
    reported separately.
  \item \textbf{Extension of the comparison to the remaining
    datasets.} If the compute budget permits, RULER, Loong, and
    LongBench-v2 will be added to both tracks.
  \item \textbf{Extension of the comparison to other
    architectures.} The remaining baselines of the program (RAPTOR,
    HippoRAG, LightRAG, PathRAG, and possibly others) will be
    executed under the same protocol.
\end{enumerate}

% =====================================================================
\section{Conclusion}

We have compared LongRAG and My\_RAG under a fully controlled,
retrieval-only protocol on four heterogeneous long-context corpora.
Under this protocol, My\_RAG achieves a macro document Recall@10 of
\num{0.4049} against LongRAG's \num{0.3144}, a relative improvement
of \num{28.79}\%. My\_RAG wins on every dataset and on every one of
the eighteen metrics reported. The magnitude of the improvement
varies systematically with the mean source-document size of the
corpus, and the largest gap appears on the corpus with the shortest
source documents. This pattern is consistent with the architectural
observation that fine-grained retrieval units preserve document-level
ranking quality on corpora whose documents are smaller than the
unit budgets of coarser architectures.

The full metric table, the controlled adaptations of both tracks,
the threats to validity, and the reproducibility chain are documented
in this report. The comparison serves as the reference result for
the next phases of the RAG Architecture Readiness Program.

% =====================================================================
\appendix

\section{Appendix A: Full Per-Dataset Metrics}

\begin{table}[h]
\centering
\caption{LongRAG per-dataset metrics.}
\small
\begin{tabular}{l r r r r r r r}
\toprule
\textbf{Dataset} & \textbf{H@1} & \textbf{H@5} & \textbf{H@10}
  & \textbf{R@10} & \textbf{nDCG@10} & \textbf{MRR@10} & \textbf{P@10} \\
\midrule
CLERC              & \num{0.0310} & \num{0.0550} & \num{0.0940}
                   & \num{0.0786} & \num{0.0543} & \num{0.0444} & \num{0.0270} \\
CUAD               & \num{0.0039} & \num{0.0161} & \num{0.0324}
                   & \num{0.0324} & \num{0.0153} & \num{0.0102} & \num{0.0032} \\
GreekBarBench      & \num{0.0352} & \num{0.1303} & \num{0.2606}
                   & \num{0.2606} & \num{0.1233} & \num{0.0828} & \num{0.0261} \\
LaRA               & \num{0.7128} & \num{0.8409} & \num{0.8861}
                   & \num{0.8861} & \num{0.7977} & \num{0.7697} & \num{0.0886} \\
\midrule
\textbf{Macro}     & \num{0.1957} & \num{0.2606} & \num{0.3183}
                   & \num{0.3144} & \num{0.2476} & \num{0.2268} & \num{0.0362} \\
\bottomrule
\end{tabular}
\end{table}

\begin{table}[h]
\centering
\caption{My\_RAG per-dataset metrics.}
\small
\begin{tabular}{l r r r r r r r}
\toprule
\textbf{Dataset} & \textbf{H@1} & \textbf{H@5} & \textbf{H@10}
  & \textbf{R@10} & \textbf{nDCG@10} & \textbf{MRR@10} & \textbf{P@10} \\
\midrule
CLERC              & \num{0.2450} & \num{0.4510} & \num{0.5430}
                   & \num{0.3546} & \num{0.2875} & \num{0.3336} & \num{0.0964} \\
CUAD               & \num{0.0046} & \num{0.0224} & \num{0.0431}
                   & \num{0.0431} & \num{0.0199} & \num{0.0130} & \num{0.0043} \\
GreekBarBench      & \num{0.0211} & \num{0.1585} & \num{0.2852}
                   & \num{0.2852} & \num{0.1314} & \num{0.0850} & \num{0.0285} \\
LaRA               & \num{0.7506} & \num{0.9067} & \num{0.9368}
                   & \num{0.9368} & \num{0.8466} & \num{0.8174} & \num{0.0937} \\
\midrule
\textbf{Macro}     & \num{0.2553} & \num{0.3846} & \num{0.4520}
                   & \num{0.4049} & \num{0.3214} & \num{0.3123} & \num{0.0557} \\
\bottomrule
\end{tabular}
\end{table}

\section{Appendix B: Retrieval Performance}

\begin{center}
\small
\begin{tabular}{l r r r}
\toprule
\textbf{Stage} & \textbf{LongRAG} & \textbf{My\_RAG} & \textbf{Hardware} \\
\midrule
Chunking / unit materialization & \SI{180}{\second} & \SI{150}{\second} & CPU \\
Document embedding & $\approx \SI{1600}{\second}$ & $\approx \SI{780}{\second}$ & T4 GPU \\
Query embedding & $\approx \SI{240}{\second}$ & $\approx \SI{240}{\second}$ & T4 GPU \\
Storage ingest & \SI{30}{\second} & \SI{960}{\second} & CPU \\
Full retrieval & \SI{342}{\second} & \SI{132}{\second} & CPU \\
Evaluation & \SI{5}{\second} & \SI{6}{\second} & CPU \\
\bottomrule
\end{tabular}
\end{center}

Both retrieval stages are comparable in wall-clock time despite the
eighteen-fold difference in corpus size, because My\_RAG's numpy
retrieval implementation is dramatically faster per query than the
Qdrant-local brute-force scan used in the earlier My\_RAG retrieval
prototype.

\section{Appendix C: Reproducibility Chain}

\begin{center}
\small
\begin{tabular}{l l}
\toprule
\textbf{Artifact} & \textbf{Location} \\
\midrule
LongRAG closure manifest & \texttt{baseline\_audit/longrag\_full\_track\_closure\_manifest.yaml} \\
My\_RAG closure manifest & \texttt{baseline\_audit/my\_rag\_m0\_closure\_manifest.yaml} \\
Cross-dataset comparison & \texttt{baseline\_audit/longrag\_vs\_my\_rag\_comparison.yaml} \\
Comparison script & \texttt{scripts/compare\_longrag\_vs\_my\_rag.py} \\
Comparison markdown & \texttt{docs/longrag\_vs\_my\_rag\_comparison.md} \\
LongRAG results & \texttt{outputs/external\_longrag/<ds>/retrieval-v1/} \\
My\_RAG results & \texttt{outputs/external\_my\_rag/<ds>/retrieval-v1/} \\
LongRAG metrics & \texttt{outputs/external\_longrag/<ds>/evaluation-v1/} \\
My\_RAG metrics & \texttt{outputs/external\_my\_rag/<ds>/evaluation-v1/} \\
\bottomrule
\end{tabular}
\end{center}

\end{document}